\section{Simulation notes: grid evaluation, cutoffs, and action rates}\label{sec:sim_notes}

\subsection{Purpose and role of the simulation}\label{subsec:sim_purpose}

The analytical results in these notes characterize best responses in terms of model primitives and
comparisons of the form
\[
    \Delta(v;S,\delta) \equiv u_1(v;S,\delta) - u_0(v;S)
    \quad\text{and}\quad
    a^\star(v;S,\delta)=\mathbf{1}\{\Delta(v;S,\delta)\ge 0\}.
\]
The work in the \texttt{computational/} folder plays a supporting role by numerically evaluating these objects on a dense grid of parameter values in order to (i) visualize
comparative statics, (ii) compute summary objects such as cutoffs and implied action rates, and
(iii) diagnose whether the best response behaves like a single-crossing cutoff in $v$ for each
$(S,\delta)$ cell.

The scripts generate \emph{model-implied} outcomes
from the primitives and the best-response rule, and then aggregate those outcomes into compact
tables that can be inspected, plotted, or used as inputs to additional analysis.

\subsection{How the model is simulated (core model \& helpers)}\label{subsec:sim_how}

All simulation code is implemented in R and is organized into:
\begin{itemize}
    \item {Core model functions in \texttt{computational/R/functions/model.R}. These define
    the payoff objects $u_0(\cdot)$ and $u_1(\cdot)$ and the incentive difference
    $\Delta(\cdot)$ (exposed as \texttt{Delta\_incentive()} in code), which is the key object used
    to construct best responses.
    \item Grid construction helpers in \texttt{computational/R/functions/grids.R}. These
    provide convenience functions to create evaluation grids over $v$ and over the parameter pair
    $(\log S,\delta)$.
    \item Manifest / reproducibility helpers in \texttt{computational/R/functions/manifest.R}.
    The main grid runner writes a small JSON manifest recording the script path, parameter settings,
    and hashes/sizes of generated outputs.
\end{itemize}

Fix $\sigma>0$ and interpret candidate valence as distributed
\[
    v \sim \mathcal{N}(0,\sigma^2).
\]
For aggregation and summary statistics, I treat candidate valence as a random draw
$v\sim\mathcal{N}(0,\sigma^2)$ for primarily computational reasons. The model objects can
be evaluated for very large $|v|$, but on a finite grid the extreme tails are both (i) empirically
less relevant and (ii) a source of numerical instability when computing ratios and
differences of probabilities. Using a normal reference distribution focuses attention on the
region of $v$ where most mass lies and yields stable numerical integrals for implied action rates.

I keep $\sigma$ explicit as a scale parameter for valence. Numerical grids and penalty magnitudes
are chosen in multiples of $\sigma$ (e.g.\ $v\in[-4\sigma,4\sigma]$), which makes the computations
comparable across alternative calibrations of the dispersion of $v$.

I grid over $\log S$ rather than $S$ directly. This ensures $S>0$ by construction and provides a
numerically well-behaved grid that samples $S$ multiplicatively (i.e.\ in equal log steps) across
orders of magnitude.

\paragraph{Step 1: evaluate incentives and best responses on a grid.}
The runner \texttt{computational/R/runners/01\_single\_agent\_grid.R} constructs:
\begin{enumerate}
    \item a dense grid for $v$ over a bounded interval (in practice, a symmetric range such as
    $[-4\sigma,4\sigma]$), and
    \item a rectangular grid over $(\log S,\delta)$ covering an economically relevant range.
\end{enumerate}
For each $(\log S,\delta)$ cell, the script computes $\Delta(v;S,\delta)$ over the entire $v$ grid and
stores the implied best response
\[
    a^\star(v;S,\delta)=\mathbf{1}\{\Delta(v;S,\delta)\ge 0\},
\]
using the same tie-break convention in all downstream scripts. The output is written as a large dataset \texttt{single\_agent\_grid.csv} under
\texttt{computational/output/datasets/}.

\paragraph{Step 2: compress to cutoffs and action rates.}
The runner \texttt{computational/R/runners/02\_cutoffs\_and\_action\_rates.R} reads the raw grid dataset and
produces a compact summary table indexed by $(\log S,\delta)$ that contains:
\begin{itemize}
    \item a cutoff $v^\star(\log S,\delta)$ defined as the \emph{first} $v$ on the grid where
    $\Delta(v;S,\delta)\ge 0$ (and $v^\star=+\infty$ if no such $v$ exists), and
    \item an implied action rate $\Pr(a^\star=1\mid \log S,\delta)$ obtained by numerically integrating
    the indicator $a^\star(v;S,\delta)$ under $v\sim\mathcal{N}(0,\sigma^2)$.
\end{itemize}
The script also diagnoses whether $a^\star(v;S,\delta)$ is consistent with a \emph{single cutoff} in $v$
(by counting the number of switches in $a^\star$ along the ordered $v$ grid). When the cutoff property holds,
the action rate can equivalently be computed from the cutoff via the normal CDF; when it fails, the script
reports the model-implied action rate computed from the full (possibly non-monotone) pattern of $a^\star(v)$.
The summary output is written to \texttt{cutoffs\_and\_action\_rates.csv} under
\texttt{computational/output/derived/}.

\paragraph{Step 3: diagnostics for non-monotone cells.}
The analysis script \texttt{computational/R/analysis/03\_diagnose\_nonmonotone\_cells.R} provides additional
checks for $(\log S,\delta)$ cells where $a^\star(v)$ switches two or more times as $v$ increases. It performs
robustness checks to small numerical tolerances, optionally recomputes $\Delta(v)$ on a much denser $v$ grid for
selected cells, and writes a compact diagnostics table plus simple plots under
\texttt{computational/output/diagnostics/}.

\subsection{Outputs and file locations}\label{subsec:sim_outputs}

The simulation pipeline produces the following:
\begin{itemize}
    \item \texttt{computational/output/datasets/single\_agent\_grid.csv}: raw grid evaluations of
    $(u_0,u_1,\Delta,a^\star)$ over $(v,\log S,\delta)$.
    \item \texttt{computational/output/derived/cutoffs\_and\_action\_rates.csv}: compressed per-cell summary
    including cutoffs, implied action rates, and monotonicity diagnostics.
    \item \texttt{computational/output/diagnostics/nonmonotone\_cells\_diagnostics.csv} and
    \texttt{computational/output/diagnostics/figures/}: targeted diagnostics for multi-switch cells.
    \item \texttt{computational/output/manifests/*.json}: manifests that record parameter settings and hashes
    for generated outputs (useful for reproducibility).
\end{itemize}

\subsection{Simulation results}\label{subsec:sim_interpret}

The simulation outputs are most useful as a diagnostic for the shape of incentives and best
responses across $(\log S,\delta)$, and for understanding when the analytical cutoff logic provides
a good approximation to global behavior.

By construction, the best response is computed from
\[
    \Delta(v;S,\delta) = u_1(v;S,\delta) - u_0(v;S)
    = p_1(v;S)\,\Phi(v-\delta) - p_0(v;S)\,\Phi(v),
\]
with $p_0(v;S)=\exp(v)/(\exp(v)+S)$ and $p_1(v;S)=\exp(v+1)/(\exp(v+1)+S)$.
Every pattern in the simulated best response can be interpreted as a contest between:
(i) a \emph{primary-stage} nomination gain (increasing $p$ via shifting $v\mapsto v+1$) and
(ii) a \emph{general-election} electability cost (shifting $\Phi(v)$ to $\Phi(v-\delta)$).

A convenient way to read the results is through the ratio condition emphasized in the analytical
notes: $a^\star=1$ is favored when the proportional nomination gain outweighs the proportional
electability cost.

The derived file \texttt{cutoffs\_and\_action\_rates.csv} summarizes each $(\log S,\delta)$ cell by a cutoff
$v^\star(\log S,\delta)$ (the first grid point where $\Delta\ge 0$) and by the implied action rate
$\Pr(a^\star=1\mid \log S,\delta)$ under $v\sim\mathcal{N}(0,\sigma^2)$.

Interpreting these objects:
\begin{itemize}
    \item When $v^\star(\log S,\delta)$ is finite and the cutoff diagnostic indicates a single switch,
    the simulation supports a clean threshold interpretation: action $a=1$ is optimal for sufficiently
    high (or sufficiently low) $v$ depending on the sign pattern of $\Delta(v)$, and the implied action
    rate is tightly linked to the normal tail probability beyond the cutoff.
    \item When $v^\star(\log S,\delta)=+\infty$ (``always 0'' in the summary), it indicates that on the
    relevant grid region, the electability cost dominates: $\Delta(v;S,\delta)<0$ for all $v$ considered.
    \item When the summary flags non-cutoff behavior (two or more switches in $a^\star(v)$), the single-crossing
    cutoff logic is \emph{not globally valid} for that parameter cell. In those regions, it is more appropriate
    to think of best responses as potentially featuring ``islands'' of $a=1$ in $v$, and to treat any cutoff-based
    analytical claims as local or conditional.
\end{itemize}

Cells with multiple switches are those where the gain--cost tradeoff is nontrivial in $v$.
the primary-stage gain from shifting $v$ by $+1$ is inherently nonlinear because $p(v;S)$ is logistic,
and the general-election term $\Phi(v-\delta)$ is also nonlinear in $v$ (especially in the tails).

Practically, this should shape how the analytical results are used:
\begin{itemize}
    \item Treat results that rely on monotonicity / single-crossing as \emph{parameter-restricted}
    statements, and use the simulation to map the region of $(\log S,\delta)$ where those restrictions
    appear satisfied.
    \item In regions with non-monotone $a^\star(v)$, summarize behavior using action rates (integrated over $v$)
    rather than a single cutoff, and interpret comparative statics in terms of changes in these rates.
\end{itemize}

The implied action rate $\Pr(a^\star=1\mid \log S,\delta)$ should be interpreted as a model-implied frequency
of choosing $a=1$ in a population of candidates with heterogeneous $v$. This object is useful because it
(i) focuses on the ``relevant'' region of $v$ under the chosen reference distribution and
(ii) remains well-defined even when $a^\star(v)$ is not a clean cutoff.

Accordingly, the action-rate surface over $(\log S,\delta)$ is a natural summary of how the model's
behavior changes with opponent strength and the severity of the general-election penalty, and it provides
a bridge between the analytical conditions and quantitative magnitudes.
